\documentclass{article}

\usepackage{amsmath,amssymb}
\usepackage{xurl}
\usepackage[hidelinks]{hyperref}
\setlength{\emergencystretch}{2em}

% Shared commands for notes in docs/paper-gaps/.

\DeclareUnicodeCharacter{2080}{\ifmmode{}_{0}\else\textsubscript{0}\fi}
\DeclareUnicodeCharacter{2081}{\ifmmode{}_{1}\else\textsubscript{1}\fi}
\DeclareUnicodeCharacter{2082}{\ifmmode{}_{2}\else\textsubscript{2}\fi}
\DeclareUnicodeCharacter{2099}{\ifmmode{}_{n}\else\textsubscript{n}\fi}

\newcommand{\Lean}{\textsc{Lean}}
\ExplSyntaxOn
\NewDocumentCommand{\leanid}{m}{
  \ifmmode
    \textsf{\detokenize{#1}}
  \else
    \group_begin:
    \urlstyle{sf}
    \tl_set:Nn \l_tmpa_tl {#1}
    \tl_replace_all:Nnn \l_tmpa_tl {\_} {_}
    \exp_args:NV \path \l_tmpa_tl
    \group_end:
  \fi
}
\ExplSyntaxOff
\providecommand{\tr}{\operatorname{tr}}
\newcommand{\tnleanIssueBase}{https://github.com/LionSR/TNLean/issues/}
\newcommand{\tnleanPrBase}{https://github.com/LionSR/TNLean/pull/}
\newcommand{\tnleanDocBase}{https://sirui-lu.com/TNLean/docs/}
\newcommand{\ghissue}[1]{\href{\tnleanIssueBase#1}{GitHub issue \##1}}
\newcommand{\ghpr}[1]{\href{\tnleanPrBase#1}{PR \##1}}
\newcommand{\doclink}[1]{\href{\tnleanDocBase#1.html}{\path{#1}}}

% Dirac notation and the identity operator, as in blueprint/src/macros/common.tex.
% \providecommand keeps note-local definitions authoritative.
\usepackage{dsfont}
\providecommand{\ket}[1]{|#1\rangle}
\providecommand{\bra}[1]{\langle #1|}
\providecommand{\braket}[2]{\langle #1 | #2 \rangle}
\providecommand{\ketbra}[2]{|#1\rangle\!\langle #2|}
\providecommand{\Id}{\mathds{1}}
\providecommand{\one}{\mathds{1}}
\providecommand{\C}{\mathbb{C}}
\providecommand{\MN}[1]{M_{#1}(\C)}
\providecommand{\MD}{M_D(\C)}

% External referencing for paper sources.  The build helper writes sanitized
% auxiliary files under build/paper-aux/ and excludes duplicated source labels.
\usepackage{xr}

% Import a generated paper auxiliary file with a caller-supplied label prefix.
% The candidate paths support builds from docs/paper-gaps/, the repository root,
% and build/paper-aux/ itself.
\newcommand{\importpaperaux}[2]{%
  \IfFileExists{../../build/paper-aux/#2.aux}{%
    \externaldocument[#1]{../../build/paper-aux/#2}%
  }{%
    \IfFileExists{build/paper-aux/#2.aux}{%
      \externaldocument[#1]{build/paper-aux/#2}%
    }{%
      \IfFileExists{#2.aux}{%
        \externaldocument[#1]{#2}%
      }{}%
    }%
  }%
}

% General external reference macros
\newcommand{\paperref}[2]{\ref{#1-#2}}
\newcommand{\papereqref}[2]{(\ref{#1-#2})}

% CPSV16 source (1606.00608 / MPDO-22-12-17-2.tex) external document and shortcuts
\importpaperaux{cpsv16-}{1606.00608/MPDO-22-12-17-2-tnlean}
\newcommand{\cpsvref}[1]{\paperref{cpsv16}{#1}}
\newcommand{\cpsveqref}[1]{\papereqref{cpsv16}{#1}}

% Verdict marker: \gapnote{<kind>}{<status>}. Kind is the policy
% classification (clarification, local-correction, scope-restriction,
% unfaithful, false-source, open-gap); status is open, wip (elimination
% underway), resolved, or historical. Machine-read by the site index;
% prints a verdict line.
\newcommand{\gapnote}[2]{\par\noindent\textbf{Verdict:} #1 (#2).\par}


\title{Correlator Exponential Decay: Rate Above the Complementary Spectral Radius}
\author{}
\date{2026-07-29}

\begin{document}
\maketitle

\section{The gap}

The source~\cite[Section~II.B.3]{Cirac2021Matrix} states the connected
correlator decays as $|C(X,Y;n)|\le C_{XY}|\lambda_2|^n$, where
$\lambda_2$ is the largest subleading eigenvalue of the transfer map
$\E_A$.  This is the exact subleading spectral radius.

The Lean formalization \leanid{MPSTensor.correlation_length_bound}
(Theorem~\ref{thm:correlation_length_bound}) obtains exponential decay
at a rate $\xi>0$ via Gelfand's formula applied to
$\E_A-P$, the complementary part of the transfer map.  Gelfand's
formula yields $\|(\E_A-P)^n\|\le Cr^n$ for any $r$ with
$\rho(\E_A-P)<r<1$, but not for $r=\rho(\E_A-P)$ itself when the
spectral radius is achieved by a non-diagonalizable Jordan block.
Consequently the rate delivered by the Gelfand route is some value
strictly above the complementary spectral radius, not the exact
subleading eigenvalue modulus.

The unconditional correlator bound proved in this issue
(\leanid{MPSTensor.connectedCorrelator_exp_decay}) therefore carries a
\textbf{scope restriction}: the rate is strictly above the
complementary spectral radius, not exactly $|\lambda_2|$.

\section{The exact rate}

The exact rate $|\lambda_2|^n$ requires an eigendecomposition of
$\E_A$ (or, more generally, a Jordan decomposition) extracting the
coefficients $c_j$ and eigenvalues $\lambda_j$.  That extraction is
the sum-of-exponentials half of Issue~\#1447 and is not delivered by
the present lemmas, which use only the norm-geometric estimate from
Gelfand's formula.

\section{What is not a discrepancy}

The softened rate is not a paper error.  The paper's derivation of the
exact rate $|\lambda_2|$ assumes a spectral expansion
$C(X,Y;n)=\sum_j c_{XY}(j)\lambda_j^n$, which the formalization has
not yet extracted.  The present bound is mathematically correct as a
consequence of the complementary transfer-map gap alone, without any
diagonalizability hypothesis.  It is stronger than the source statement
in the sense that it does not require the transfer map to be
diagonalizable, and weaker in that the rate is not exactly the
subleading spectral radius.

\bibliographystyle{alpha}
\bibliography{references}

\end{document}
